% part: intuitionistic-logic
% chapter: semantics
% section: topological-semantics

\documentclass[../../../include/open-logic-section]{subfiles}

\begin{document}

\olfileid{int}{sem}{top}

\olsection{معناشناسی توپولوژیک}

راه دیگری برای ارائهٔ معناشناسی منطق شهودگرایانه استفاده از مفهوم
ریاضی توپولوژی است.

\begin{defn}
  فرض کنید~$X$ یک مجموعه باشد. یک \emph{توپولوژی روی~$X$} مجموعه‌ای
  $\Top{O} \subseteq \Pow{X}$ است که خواص زیر را برآورده می‌کند.
  !!{element}های~$\Top{O}$ \emph{مجموعه‌های باز} توپولوژی نامیده
  می‌شوند. مجموعهٔ~$X$ به‌همراه~$\Top{O}$ یک \emph{فضای توپولوژیک}
  نامیده می‌شود.
  \begin{enumerate}
  \item مجموعهٔ تهی و کل فضا بازند: $\emptyset$، $X \in \Top{O}$.
  \item مجموعه‌های باز تحت اشتراک‌های متناهی بسته‌اند: اگر~$U$،
    $V \in \Top{O}$، آنگاه~$U \cap V \in \Top{O}$.
  \item مجموعه‌های باز تحت اجتماع‌های دلخواه بسته‌اند: اگر برای هر
    $i \in I$، $U_i \in \Top{O}$، آنگاه
    $\bigcup \Setabs{U_i}{i \in I} \in \Top{O}$.
  \end{enumerate}
\end{defn}

اگر مجموعهٔ مجموعه‌های باز از بافت معلوم باشد، ممکن است برای فضای
توپولوژیک صرفاً~$X$ بنویسیم؛ اما توجه کنید که~$X$ تنها پس از مجهزشدن
به توپولوژی~$\Top{O}$ یک فضای توپولوژیک نامیده می‌شود.

\begin{defn}
  یک \emph{مدل توپولوژیک} برای منطق گزاره‌ای شهودگرایانه سه‌تاییِ
  $\mModel{X} = \tuple{X, \Top{O}, V}$ است، که در آن~$\Top{O}$ یک
  توپولوژی روی~$X$ و~$V$ تابعی است که به هر متغیر گزاره‌ای یک مجموعهٔ
  باز در~$\Top{O}$ نسبت می‌دهد.

  با داشتن مدل توپولوژیک~$\mModel{X}$، می‌توانیم~$\Prop{X}{!A}$ را
  به‌طور استقرایی چنین تعریف کنیم:
  \begin{enumerate}
  \item $\Prop{X}{\lfalse} = \emptyset$
  \item $\Prop{X}{p} = V(p)$
  \item $\Prop{X}{!A \land !B} = \Prop{X}{!A} \cap \Prop{X}{!B}$
  \item $\Prop{X}{!A \lor !B} = \Prop{X}{!A} \cup \Prop{X}{!B}$
  \item $\Prop{X}{!A \lif !B} = \Interior{(X \setminus \Prop{X}{!A}) \cup
    \Prop{X}{!B}}$
  \end{enumerate}
  در اینجا~$\Interior{V}$ تابعی است که مجموعهٔ~$V \subseteq X$ را به
  \emph{درون} آن، یعنی اجتماع همهٔ مجموعه‌های بازی که در خود دارد،
  می‌فرستد. به بیان دیگر،
  \[
  \Interior{V} = \bigcup \Setabs{U}{U \subseteq V \text{ و } U \in \Top{O}}.
  \]
\end{defn}

توجه کنید که درون هر مجموعه همواره باز است، زیرا اجتماعی از مجموعه‌های
باز است. بنابراین~$\Prop{X}{!A}$ همواره مجموعه‌ای باز است.

هرچند معناشناسی توپولوژیک بسیار انتزاعی است، می‌توان آن را به شیوه‌هایی
فهمید که انگیزه‌اش را روشن کنند. فرض کنید !!{element}های~$X$، یا
«نقطه‌های» آن، نقاطی باشند که گزاره‌ها در آن‌ها قابل ارزیابی‌اند.
مجموعهٔ همهٔ نقاطی که~$!A$ در آن‌ها صادق است، گزارهٔ بیان‌شده با~$!A$
است. هر مجموعه‌ای از نقاط یک گزارهٔ ممکن نیست؛ تنها !!{element}های
$\Top{O}$ چنین‌اند. $!A \Entails !B$ اگر و تنها اگر در هر~$X$، $!B$
در هر نقطه‌ای که~$!A$ صادق است صادق باشد؛ یعنی
$\Prop{X}{!A} \subseteq \Prop{X}{!B}$. گزارهٔ محال~$\lfalse$ هرگز
صادق نیست، پس~$\Prop{X}{\lfalse} = \emptyset$.

برای آنکه قوانین معتبر شهودگرایانه برقرار باشند، گزاره‌های بیان‌شده با
$!B \land !C$، $!B \lor !C$ و~$!B \lif !C$ باید چه نسبتی با گزاره‌های
بیان‌شده با~$!B$ و~$!C$ داشته باشند؟ یعنی چگونه داشته باشیم
$!A \Proves !B$ اگر و تنها اگر
$\Prop{X}{!A} \subseteq \Prop{X}{!B}$؟ برای هر~$!A$،
$\lfalse \Proves !A$ را لازم داریم؛ و چون برای هر~$U$،
$\emptyset \subseteq U$، این شرط برقرار است. چون
$!B \land !C \Proves !B$، لازم داریم
$\Prop{X}{!B \land !C} \subseteq \Prop{X}{!B}$؛ و به همین ترتیب
$\Prop{X}{!B \land !C} \subseteq \Prop{X}{!C}$. بزرگ‌ترین مجموعه‌ای
که~$W \subseteq U$ و~$W \subseteq V$ را برآورده می‌کند~$U \cap V$
است. برعکس، $!B \Proves !B \lor !C$ و~$!C \Proves !B \lor !C$؛ پس
لازم داریم~$\Prop{X}{!B} \subseteq \Prop{X}{!B \lor !C}$ و
$\Prop{X}{!C} \subseteq \Prop{X}{!B \lor !C}$. کوچک‌ترین مجموعهٔ~$W$
که~$U \subseteq W$ و~$V \subseteq W$ را برآورده می‌کند~$U \cup V$
است.

تعریف~$\lif$ ظریف است: $!A \lif !B$ ضعیف‌ترین گزاره‌ای را بیان می‌کند
که همراه با~$!A$، $!B$ را نتیجه می‌دهد. از
$(!A \lif !B) \land !A \Proves !B$ روشن است که ترکیب~$!A \lif !B$
با~$!A$، $!B$ را نتیجه می‌دهد. پس~$\Prop{X}{!A \lif !B}$ باید
بزرگ‌ترین مجموعهٔ بازی باشد که
$\Prop{X}{!A \lif !B} \cap \Prop{X}{!A}
\subseteq \Prop{X}{!B}$؛ و این به تعریف ما می‌انجامد.

\end{document}
